\documentclass[a4paper]{article}

\usepackage{lmodern}
\usepackage[margin=25mm]{geometry}
\usepackage[hidelinks]{hyperref}
\usepackage{array}
\usepackage{graphicx}
\usepackage{verbatim}
\usepackage{corasdiagram}

\title{\texttt{corasdiagram} Manual}
\date{Version \corasdiagramversion}

\newcommand{\exampletex}[1]{../examples/#1.tex}
\newcommand{\examplepdf}[1]{../examples/#1.pdf}
\newcommand{\examplebox}[1]{%
  \fbox{\parbox{.9\linewidth}{Build \texttt{examples/#1.tex} first to embed the rendered output here.}}%
}
\newcommand{\completeexample}[2]{%
  \subsection{#2}
  \paragraph{Compileable source}
  \verbatiminput{\exampletex{#1}}
  \paragraph{Rendered output}
  \begin{center}
  \IfFileExists{\examplepdf{#1}}{%
    \includegraphics[width=.96\linewidth]{\examplepdf{#1}}%
  }{%
    \examplebox{#1}%
  }
  \end{center}
}

\begin{document}

\maketitle
\tableofcontents

\section{Overview}

\texttt{corasdiagram} is a semantic-first LaTeX package for CORAS diagrams and
tables. The preferred public API uses semantic nouns and verbs that mirror the
formal English-prose interpretation of the CORAS language.

The canonical diagram environments are:

\begin{itemize}
  \item \texttt{corasassetdiagram}
  \item \texttt{corasthreatdiagram}
  \item \texttt{corasriskdiagram}
  \item \texttt{corastreatmentdiagram}
\end{itemize}

Treatment overview diagrams are expressed with
\texttt{corastreatmentdiagram[mode=overview]}. The package also ships the
high-level analysis and high-level risk tables.

This manual is scoped to printable reference use:

\begin{itemize}
  \item semantic-first commands and keys
  \item compact compatibility appendix
  \item advanced structural appendix
  \item complete compileable examples of every diagram and table family
\end{itemize}

For the user-facing migration guide, see
\texttt{docs/semantic-api-migration.md}. Maintainer-oriented release and API
inventory material lives outside the public documentation set.

\section{Installation}

Compile from the repository root with:

\begin{verbatim}
TEXINPUTS=tex/latex//: latexmk -pdf -interaction=nonstopmode -halt-on-error \
  examples/corasdiagram-minimal.tex
\end{verbatim}

For a user TEXMF installation, copy \texttt{tex/latex/corasdiagram/} into
\texttt{\$HOME/texmf/tex/latex/corasdiagram/}.

\section{Semantic-First DSL}

\subsection{Asset diagrams}

English sentence forms:

\begin{itemize}
  \item ``... is a party.''
  \item ``... is a direct asset.''
  \item ``... is an indirect asset.''
  \item ``... assigns value to ....''
  \item ``... assigns the value ... to ....''
  \item ``Harm to ... may result in harm to ....''
\end{itemize}

Canonical commands:

\begin{itemize}
  \item \texttt{\textbackslash party[...]}
  \item \texttt{\textbackslash asset(direct)[...]}
  \item \texttt{\textbackslash asset(indirect)[...]}
  \item \texttt{\textbackslash assigns\{party -> asset\}[asset value=...]}
  \item \texttt{\textbackslash harms\{asset -> asset\}}
\end{itemize}

\subsection{Threat diagrams}

English sentence forms:

\begin{itemize}
  \item ``... is a deliberate human threat.''
  \item ``... is an accidental human threat.''
  \item ``... is a non-human threat.''
  \item ``... is a vulnerability.''
  \item ``Threat scenario ... occurs with likelihood ....''
  \item ``Unwanted incident ... occurs with likelihood ....''
  \item ``... exploits vulnerability ... to initiate ....''
  \item ``... leads to ... with conditional likelihood ....''
  \item ``... impacts ... with consequence ....''
\end{itemize}

Canonical commands:

\begin{itemize}
  \item \texttt{\textbackslash threat(deliberate|accidental|nonhuman)[...]}
  \item \texttt{\textbackslash vulnerability[...]}
  \item \texttt{\textbackslash threatscenario[...]}
  \item \texttt{\textbackslash unwantedincident[...]}
  \item \texttt{\textbackslash initiates\{source -> target\}[vulnerability=...]}
  \item \texttt{\textbackslash leadsto\{source -> target\}[conditional likelihood=...]}
  \item \texttt{\textbackslash impacts\{incident -> asset\}[consequence=...]}
\end{itemize}

\subsection{Risk diagrams}

English sentence forms:

\begin{itemize}
  \item ``Risk ... occurs with risk level ....''
  \item ``... initiates ....''
  \item ``... impacts ....''
\end{itemize}

Canonical commands:

\begin{itemize}
  \item \texttt{\textbackslash risk[...]}
  \item \texttt{\textbackslash initiates\{source -> risk\}}
  \item \texttt{\textbackslash impacts\{risk -> asset\}}
\end{itemize}

\subsection{Treatment diagrams}

English sentence forms:

\begin{itemize}
  \item ``... is a treatment scenario.''
  \item ``... treats ....''
  \item ``... decreases the likelihood of ....''
  \item ``... decreases the consequence of ....''
\end{itemize}

Canonical commands:

\begin{itemize}
  \item \texttt{\textbackslash treatmentscenario[...]}
  \item \texttt{\textbackslash treats\{treatmentscenario -> target\}[treatment category=...]}
  \item plus \texttt{\textbackslash initiates}, \texttt{\textbackslash leadsto}, and
    \texttt{\textbackslash impacts} when the treated threat or risk chain is shown
\end{itemize}

\subsection{Treatment overview mode}

Use \texttt{corastreatmentdiagram[mode=overview]} when treatments are attached
directly to risks instead of a full threat/risk chain.

\section{Tables}

\subsection{High-level analysis table}

\texttt{corashighlevelanalysistable} is the canonical high-level table
environment. Use \texttt{\textbackslash corashighlevelanalysisrow} to add rows.

\subsection{High-level risk table}

\texttt{corashighlevelrisktable} and
\texttt{\textbackslash corashighlevelriskrow} remain available for the risk
table naming variant. They render the same table structure as the analysis
table and are kept for compatibility and document readability.

\section{Canonical Keys}

\subsection{Node keys}

\begin{center}
\begin{tabular}{>{\ttfamily}p{0.24\linewidth}p{0.24\linewidth}p{0.38\linewidth}}
\hline
Key & Applies to & Meaning \\
\hline
id & all typed nodes & Stable public identifier. \\
risk ref & risk, unwantedincident & External reference tag such as \texttt{R4}. \\
risk level & risk & Risk level line such as \texttt{unacceptable}. \\
likelihood & threatscenario, unwantedincident & Plain-text likelihood such as \texttt{possible}. \\
likelihood basis & threatscenario, unwantedincident & Bracketed quantitative basis such as \texttt{3,15:30y}. \\
at & all typed nodes & Manual placement coordinate. \\
right / left / above / below & all typed nodes & Relative placement helpers. \\
scope & asset and party nodes & Named scope fit group membership. \\
order & auto-laid-out families & Preferred auto-layout order. \\
\hline
\end{tabular}
\end{center}

\subsection{Edge keys}

\begin{center}
\begin{tabular}{>{\ttfamily}p{0.24\linewidth}p{0.24\linewidth}p{0.38\linewidth}}
\hline
Key & Applies to & Meaning \\
\hline
asset value & \textbackslash assigns & Rendered asset-value label. \\
vulnerability & \textbackslash initiates & One or more grouped vulnerabilities on the edge. \\
conditional likelihood & \textbackslash leadsto & Conditional likelihood label on the edge. \\
consequence & \textbackslash impacts & Consequence label on the edge. \\
treatment category & \textbackslash treats & Treatment-category label on the edge. \\
pos & edge commands & Label position override. \\
path / route / tikz / label options & edge commands & Drawing overrides and styling. \\
\hline
\end{tabular}
\end{center}

\section{Compatibility Appendix}

Retained compatibility aliases:

\begin{itemize}
  \item \texttt{corastreatmentoverviewdiagram}
  \item \texttt{\textbackslash asset[...]} as \texttt{\textbackslash asset(direct)[...]}
  \item \texttt{\textbackslash indirectasset[...]}
  \item \texttt{\textbackslash scenario[...]}
  \item \texttt{\textbackslash incident[...]}
  \item \texttt{\textbackslash treatment[...]}
  \item \texttt{\textbackslash cause\{...\}}
  \item \texttt{\textbackslash impact\{...\}}
  \item \texttt{\textbackslash relate\{...\}}
  \item \texttt{\textbackslash concern\{...\}}
  \item \texttt{\textbackslash treat\{...\}}
  \item \texttt{\textbackslash associate\{...\}}
  \item \texttt{risk id}, \texttt{risk-id}, and \texttt{riskid}
\end{itemize}

Use the semantic-first commands in new documents. The migration guide covers
one-to-one replacements and the \texttt{\textbackslash cause} dispatcher rule.

\section{Advanced Structural Appendix}

The following commands remain public, but they are structural helpers rather
than part of the first-line semantic vocabulary:

\begin{itemize}
  \item \texttt{\textbackslash junction}
  \item \texttt{\textbackslash corasscope}
  \item \texttt{\textbackslash corascontainer}
  \item \texttt{\textbackslash corasriskref}
  \item \texttt{\textbackslash corasnode}
  \item \texttt{\textbackslash corasedge}
  \item \texttt{\textbackslash associate}
\end{itemize}

\section{Complete Examples}

\completeexample{corasdiagram-asset-diagram}{Asset diagram}
\completeexample{corasdiagram-threat-diagram}{Threat diagram}
\completeexample{corasdiagram-risk-diagram}{Risk diagram}
\completeexample{corasdiagram-treatment-diagram}{Treatment diagram}
\completeexample{corasdiagram-treatment-overview-diagram}{Treatment overview diagram}
\completeexample{corasdiagram-high-level-analysis-table}{High-level analysis table}
\completeexample{corasdiagram-high-level-risk-table}{High-level risk table}

\end{document}
